% !TeX root = ../papers/nc-p3.tex
\section{有限点枝の分離三次式とスキーム上の標準四移動}
\label{sec:finite-cubic-separators}
四辺が有限な標準有限Frobenius候補とし、点グラフを $X_i$ の間の同型の列と見る。
標準座標の四移動で $X_{i+4}=X_i$ とする。
\begin{theorem}\label{thm:finite-cubic-separators}
全ての有限Frobenius候補で、標準座標における四移動は点スキーム上で恒等である。
反対階数も元代数の正則性も仮定しない。
中支持3の点を一つでも持つ候補は全次数で正則であり、未認証候補では全ての中支持が4である。
任意の閉点 $p\in X_i$ を固定し、反対位置の変数 $y$ において $w(p,y)$ を
両端の三支持の間の $3\times3$ 行列 $M_p(y)$ と読む。
$\det M_p$ は非零の三次式で、反対位置の長さ20のスキーム上では二移動した対応点の
局所socleを張る。特にその点が被約なら当該点で行列の階数は3である。
\end{theorem}
\begin{proof}
前節により各二次関係空間は可分解元を含まず、中支持の次元は3または4である。
中支持4の三テンソルの両端を基底化する。行列 $M(y)$ の四係数は独立であり、
行または列を非零定数ベクトルで縮約した $3\times4$ 行列はいずれも階数2以上である。
これはそれらが隣接関係の元であり、全一因子支持と純粋関係の不在を使うためである。

$\det M$ が恒等零なら一般階数は2である。一般階数1以下の線型行列空間は
共通な像直線か共通な余像直線を持ち、端支持3に反する。
余因子行列は多項式環上の階数1の二次行列で $c(y)u(y)v(y)^t$ と因数分解できる。
実際、有理関数体上で階数1の積にし、UFD上で二ベクトルの分母と成分の最大公約因子を取り除けばよい。
$M$ は定数左右核を持たないので両ベクトルは正次数であり、従ってともに一次である。
$M(y)u(y)=0$ とする。$u$ の像が一次元なら定数核がある。
二次元なら基底を替えて $u=(y_1,y_2,0)^t$ となり、最初の二列は $y_2 c,-y_1 c$
という定数列 $c$ の倍となる。非零な一列が階数1の切断となり矛盾する。
三次元なら $u=(y_1,y_2,y_3)^t$ としたとき、第四係数行列が三つの基底ベクトルを消すので零である。
これも中支持4の係数の独立性に反する。

中支持3の場合には前節で得た三文字部分代数が標準三次元AS正則であり、
三つの縮約平面に階数1の行列はない。その三行三列の行列式が恒等零なら、
三文字の点対応は射影平面全体の同型のグラフである。
従って元の一辺の六関係のうち三つが二つの三文字空間の間の可換関係に基底変更で一致する。
ところがこの三関係の零点は、両射影三空間の三座標が比例する四次元の閉部分多様体を含む。
実際、共通の射影二座標比が決まった後で双方の補座標がそれぞれ一自由度を持つ。
この四次元多様体と残り三つのSegre超平面の交叉は次元1以上であり、有限性に反する。
従って中支持3の場合も行列式は非零である。

変数 $p$ の属する局所Artin成分を $S$ とする。$p$ から反対位置へ前向き・後向きに
二移動する射を $f,g:S\to X_{i+2}$ とする。
$M_p(y)$ の両端支持は $S$ 上の階数3の部分束となるので、局所基底を選んで
$\delta=\det M_p(y)$ を $S$ 係数の三次式とする。反対文字の全三次式の制限は
$D=H^0(X_{i+2},\mathcal O)$ への同型なので、直線束を自明化して
$\delta\in\mathcal O_S\otimes D$ と読める。閉ファイバーでこれは非零である。

$S\times X_{i+2}$ 上では $y$ 自身の先頭・末尾評価も階数3である。
従って $w(p,y)$ の各外側の支持は、$p$ により定まる超平面と $y$ により定まる超平面の
双方に含まれる。$M_p(y)$ の三列を、$p$ 側の超平面に制限した $y$ 側の点で縮約すると零になる。
余因子行列を掛けると、その点の三つの制限座標は全て $\delta$ を消す。
これら三座標が生成するのは二つの超平面が等しいという対角イデアルの引戻しである
（単元座標で規格化すれば三つの座標差である）。
片端ではこれが $f$ のグラフイデアル $I_f$、他端では $g$ のグラフイデアル $I_g$ となる。
これは隣接辺の射影がどちらも同じ点スキームへの閉埋込みであることによる。

有限Gorenstein環 $D$ の完全traceを用いると
$\operatorname{ann}(I_f)=\operatorname{Hom}_{\mathcal O_S}(\mathcal O_S,\mathcal O_S)$
に可逆層の自明化を除き同定できる。具体的には相対完全積対合のもとで
グラフ商 $\mathcal O_S$ の相対双対に等しく、その形成は閉ファイバーへの制限と可換である。
従って閉ファイバーで非零の $\delta$ は $\operatorname{ann}(I_f)$ を張る。
同じく $\operatorname{ann}(I_g)$ を張り、直交補を二度取れば $I_f=I_g$ を得る。
これで $f=g$ であり、前へ四移動すると $S$ 上で戻る。閉点に制限した消去イデアルは
対応点の極大イデアルの零化、すなわちsocleの直線である。
第一の場合の正則性は定理\ref{thm:ternary-point-subalgebras}と全四位置の純粋関係の不在で従う。
\end{proof}

\begin{corollary}\label{cor:reduced-point-six-mixed-relations}
被約な局所成分を一つ選び、その四周期列に沿って $H_i=\ker p_i$ とし、
$p_i(t_i)=1$ なる補文字を取る。各位置には二つの線型写像 $F_i,G_i$ が一意にあり、
元の六関係は
\[
 t_i v=F_i(v),\qquad u t_{i+1}=G_i(u),\qquad
 u\in H_i,\quad v\in H_{i+1},\quad F_i(v),G_i(u)\in H_i\otimes H_{i+1}
\]
となる。任意の正の長さの経路空間は、全て $t$ の一語と $H$ だけの語の像の直和である。
\end{corollary}
\begin{proof}
$p_i,p_{i+1}$ は辺をなすので $t_i t_{i+1}$ の係数は各関係で零である。
先頭が $H_i$ にある三関係を $p_{i+1}$ で縮約する写像の階数は
反対点 $p_{i-1},p_{i+1}$ での行列の階数に等しい。実際この三関係は
$w(p_{i-1})$ の末尾縮約像である。選んだ点は被約なので階数3であり、
先頭と末尾の双方が $H$ にある関係はない。従って六関係から
$t_i H_{i+1}\oplus H_i t_{i+1}$ への射影は同型である。一つでも $H$ がある語では、
その隣の補文字を順に消せる。全補文字の語の独立性は点列での評価で分かる。
\end{proof}
